\documentclass[12pt]{article}

\usepackage{coffeestains}
\usepackage{cmap}
\usepackage[T2A]{fontenc}
\usepackage[utf8]{inputenc}
\usepackage[english,russian]{babel}
\usepackage{tabularx}
\usepackage{graphicx}
\usepackage{amsmath}
\usepackage{indentfirst}
\graphicspath{ {./images/} }

\usepackage{scrextend}
\usepackage{titlesec}  % <-- добавьте этот пакет

% Переопределяем формат заголовка раздела
\titleformat{\section}[display]
  {\normalfont\huge\bfseries}        % формат всего заголовка
  {\sectionname\ \thesection}        % "Лекция 1"
  {0pt}                               % отступ между номером и темой
  {\Huge}                             % формат темы
\titlespacing*{\section}{0pt}{0pt}{20pt}  % отступы: слева, до, после

% Переопределение названия
\addto\captionsrussian{
  \renewcommand{\sectionname}{Лекция}
}

\usepackage{xcolor}
\definecolor{pastelblue}{rgb}{0.0, 0.2, 0.4}
\usepackage[colorlinks=true, linkcolor=pastelblue]{hyperref}

\title{Конспект недостающих четырёх лекций Михаила Густокашина по дисциплине «Алгоритмы и структуры данных» 4-го модуля}
\author{От Добросовестной Слушательницы Ли Людмилы БПМИИ247}
\date{2025 год}

\begin{document}
\maketitle
\newpage
\tableofcontents
\newpage

\section*{Лекция 9}
\section*{Дерево отрезков}
\addcontentsline{toc}{section}{Лекция 9: Дерево отрезков}
На merge и split из предыдущей лекции уходит довольно много времени, ведь у них большая константа. Пусть наша задача --- это находить элемент и искать минимум на отрезке. Возьмём массив, представленный на картинке. Мы для каждой пары элементов посчитаем то, что нас просят: минимум. Если кол-во элементов в массиве нечетное, то мы добавляем нейтральный по минимуму элемент --- это $-\infty$. Добавлять этот элемент нужно до тех пор, пока кол-во чисел в массиве не станет равным степени двойки.

\begin{figure}[h]
\centering
\includegraphics[width=0.56\linewidth]{gust7.png}
\caption{Дерево отрезков}
\end{figure}

На самом нижнем уровне каждый элемент соответствует отрезку длинной один. Правая и левая границы включительно. Я приписываю каждому элемента дерева какому отрезку от соответствует. Его глубина строго равна $logN$. 

\begin{figure}[h]
\centering
\includegraphics[width=0.56\linewidth]{gust8.png}
\caption{Дерево отрезков 2.0}
\end{figure}

Научимся считать минимум на отрезке. Например, от $2$ до $5$. Мы всегда начинаем считать с корня. Если очередной элемент (в нашем случае корень) полностью покрылся запросом, то мы просто выплёвываем результат верха. В нашем случае это не так. Значит, если элемент не покрывается запросом, то мы смотрим на левого и правого ребенка. И с теми детьми, которыми пересекаются запросами, для них мы, собственно, запускаем этот запрос. Такая рекурсивная функция. 

\begin{figure}[h]
\centering
\includegraphics[width=0.65\linewidth]{gust9.png}
\caption{Дети}
\end{figure}

Почему она работает именно за $logN$? Когда наша функция раздваивается по-настоящему? Когда он охватывает наш узел, то есть вот наш узел, его дети соответствуют двум половинкам этого отрезка. Приходит какой-то запрос, который покрывает серединку, и теперь наш путь раздваивается, ведь мы уходим в двух детей сразу. Но теперь запрос будет обязательно иметь границу: в левом ребенке правую границу накрывать, в правом ребенке --- левую границу. Далее у нас снова может раздвоиться, но в нашем случае в левом ребенке правый ребенок полностью покроется, поэтому мы вернемся за $O(1)$. А для всех остальных запросов они будут прижаты для всех остальных границ отрезка. Индексация у нас с нуля! Подсчитывать числа можно не явно, а динамически. Если вы не верите в то, что вы напишите нормально, то  вы можете написать с этими массивами, отладить, а потом при необходимости удалить. 

Если мы меняем какой-то элемент, то мы перестраиваем всю ветку, к которой он принадлежит. То есть мы идем рекурсией, когда локализовываем новый объект, а потом пересчитываем всех его предков. 

Как это хранить? По-честному, заводить дерево с левым и правым указателем НОДа и т.д. Но можно лучше. Смотрите, наше дерево обладает уникальным свойством: оно полное, то есть оно имеет либо двух детей, либо ни одного. В куче похожее свойство. Поэтому наше дерево можно хранить в обычном одномерном массивчике, кол-во элементов которого будет равно кол-ву элементов в дереве (не забываем про степень двойки). Нам нужно заводить массив размером $2N-1$. Наши исходные элементы мы записываем во вторую половинку массива и так далее. Массив на картинке снизу соответствует нашей куче.

\begin{figure}[h]
\centering
\includegraphics[width=0.65\linewidth]{gust100.png}
\caption{Массив}
\end{figure}

Теперь рассмотрим групповые операции. Реализуем операцию прибавления на отрезке. Как это делать? Точно так же как в декартовом дереве. Если накрывается запросом, то вы накрываете обещанием, а потом все пересчитываете. В случае, если ваш запрос проходит через ваш элемент, то по-честному выполняете обещание, этого обещания лишаетесь и ваши дети приобретают обещание и +1. Тут делается абсолютно то же самое, что и в декартовом дереве. 


\subsection*{Сжатое дерево отрезков}
\addcontentsline{toc}{subsection}{Сжатое дерево отрезков}
Теперь рассмотрим задачу. У нас есть отрезок $[0, 10^9]$. У нас есть три операции: прибавление, удаление и нахождение кол-ва точек на отрезке от $L$ до $R$. Если бы у нас было $10^5$, то у нас было бы обычное дерево отрезков. (Как считать сумму на дереве отрезков? Просто детей складываем и всё). Но $10^9$ это чересчур много. Предположим, вам лень писать декартово дерево. \textbf{Идея}: мы мысленно представим себе, что у нас есть дерево отрезков, у которого на последнем слое есть те самые элементы от $0$ до $10^9$. Но у нас нет столько памяти, поэтому мы поступим так: в тех узлах, в которых у нас ничего нету, мы их просто не задаем. Условно нам говорят: добавь тройку. Хорошо, мы ушли в левого ребенка, затем снова влево и т.д. 

\begin{figure}[h]
\centering
\includegraphics[width=0.64\linewidth]{gust101.png}
\end{figure}

В итоге у нас есть отрезок, который соответствует длине $0-3$, у него появляется ребенок $2-3$, а у него появляется ребенок $3-3$. Так это происходит. Мы \textbf{не} храним те узлы, которые соответствует отрезкам, на которых нет точек. А когда у нас появляется точка, мы все эти узлы создаём. Их будет логарифм. Потом мне скажут добавить единичку, и теперь я буду идти по существующим узлам, и к тройке создам второго ребенка $0-1$, и ему еще ребенка $1-1$ (см. рис.)

\begin{figure}[h]
\centering
\includegraphics[width=0.64\linewidth]{gust102.png}
\end{figure}

Вот здесь единственный способ хранения --- это динамическая память с указателем направо и налево. Отвечать на запрос в таком случае мы будет за $log10^9$. 

Вот есть у меня какое-то дерево отрезков, и я хочу реализовать операцию запроса на минимум. Но теперь нужно сказать не только сам минимум, но и кол-во раз его появления на отрезке. Создать нужно поле с кол-вом минимумов. Вот мы находимся в узле, значит, дети умеют говорить свой минимум и своё количество. Очевидно, что значение минимума --- это минимум из минимумов детей, а их кол-во минимумов? Если они не совпали, то мы просто берем тупо их копируем. Допустим, $min_1$ оказался меньше, чем $min_2$, поэтому мы просто скопируем. Если они совпали, то мы суммируем количество. Теперь усложним задачу: нам нужно сказать количество нулей на отрезке. 

\paragraph{Задача о камерах.} 
Предположим, у вас на участке стоят камеры и теперь нам нужно найти кол-во точек, которые не покрываются камерами. Операция установки камеры --- это просто прибавление единички на всём отрезке, то есть каждая точка этого отрезка просматривается одной камерой. С отложенным обещаем и всем на свете. Когда мы хотим узнать, сколько точек не просматривается, это количество нулей. Как найти кол-во нулей на отрезке? Воспользуемся предыдущей задачей: нам просто нужно делать запрос на минимум и кол-во минимумов на отрезке. Если минимум совпал с нулём, то ура, выводим это кол-во. Если минимум не совпал с нулём, значит, ноль нулей у нас было. В случае расстановки камер отрицательных чисел быть не может. 

\begin{figure}[h]
\centering
\includegraphics[width=0.64\linewidth]{gust104.png}
\end{figure}

Теперь нам нужно найти $k$-ый ноль. Нам нужно научиться искать индекс $k$-го нуля. А давайте бинпоиском найдём индекс этого нуля. Делаем запрос: сколько же нулей от $0$ до $i$? Если слишком много нулей, то меня правую границу поиска. Если слишком мало --- левую. Всё круто, кроме того, что у нас лишний логарифм. Сейчас мы от него избавляться. Нам нужно найти $k$-ый ноль. Левый и правый сын фактически знают кол-во нулей. Если $k$ меньше, чем количество слева, то мы уходим в левое поддерево. Если $k$ больше, чем количество слева, то мы уходим в правое поддерево и ищем там уже $k-cntz_1$ нулей (см. рис). Тут совершенно другая структура: просто элемент предок соответствует большему отрезку. 

\begin{figure}[h]
\centering
\includegraphics[width=0.5\linewidth]{gust999.png}
\end{figure}

Следующая задача: посчитать кол-во элементов, которые меньше или равны $x$ на отрезке от $L$ до $R$. Чтобы упростить задачу, не будем использовать обещания. Сделаем такую штуку: мы имеем какие-то значения, и мы в каждом узле будем хранить вообще все элементы, которые лежат на этом отрезке, в отсортированном порядке. 

\begin{figure}[h]
\centering
\includegraphics[width=0.64\linewidth]{gust107.png}
\end{figure}

Как мы это будем делать? А делать мы это будем с помощью сортировки слиянием. Внизу понятно: в листьях лежит по одному элементу. Как мы в предке узнаем, какие элементы будут в отсорченном порядке будут лежать. Мы берём отсорченную последовательность левого и правого детей и mergeим их сортировкой слиянием, выбирая просто очередное значение. А как теперь научиться отвечать на запрос? Если узел накрылся полностью, мы тупо делаем в нём бинпоиск. А если нет, то это не наши проблемы и мы запустили бинпоиск для детей, которые пересекаются с запросом. Они нам сказали ответ, а мы его просто просуммировали. Занимает это $NlogN$ памяти. А по времени построение этого всего занимает тоже $NlogN$. Если оценить количество запросов сверху, то получается логарифм квадрат. А если мы всё-таки вернем изменение одного элемента, то с массивами уже так не получится, поэтому в каждом узле дерева отрезков нам придётся хранить сбалансированное бинарное дерево поиска. После этого мы можем реализовать операцию изменения одного элемента. Когда мы меняем один элемент, мы на всём пути от корня до этого элемента убираем из multiset старое значение,и добавляем в multiset новое значение. Это тоже работает за логарифм. Единственная проблема: multiset не умеет делать определения индекса по значению. Именно поэтому вы пишете ручками сбалансированное бинарное дерево поиска, которое умеет делать эту операцию за логарифм. 

\subsection*{Многомерное дерево отрезков}
\addcontentsline{toc}{subsection}{Многомерное дерево отрезков}
Задача простая: нужно как с точечками. Операции: поставь точечку, удали точечку и посчитай кол-во точечек в прямоугольнике. Приходит запрос в виде прямоугольника и нужно подсчитать кол-во точечек, попадающих в этот прямоугольник. Чтобы решить задачу, мы будем делать дерево отрезков по одной координате, например, по $x$. Вот у нас есть половинки, и внутри каждого узла лежит дерево уже по $y$. И этот $y$ соответствует $x$. В каждом узле дерева отрезков лежит дерево отрезков. Как мы отвечаем на запросы? Мы решаем задачу по $x$, и каждый раз, когда отрезок накрылся полностью, для этого отрезочка решаем задачу по $y$. А уже там решаем задачу так, как мы уже умеем. Сложность: логарифм квадрат. Как добавлять точку? Мы смотрим все деревья отрезков по $x$, в каждом из них лежит дерево отрезков по $y$, и мы в это дерево отрезков по $y$ добавляем нашу точку. Сложность: логарифм квадрат. Это жрёт много памяти. Тогда мы пишем сжатое дерево отрезков.

\begin{figure}[h]
\centering
\includegraphics[width=0.7\linewidth]{gust108.png}
\end{figure}





































\newpage
\section*{Лекция 13}
\section*{Комбинаторные объекты}
\addcontentsline{toc}{section}{Лекция 13: Комбинаторные объекты}
Обсудим способы оптимизации перебора. Существует большое кол-во задач, которое не решается полиномиальными алгоритмами. Те алгоритмы, которые мы с вами реализовывали, как раз таки являются полиномиальными. То есть существует некий многочлен, который описывает кол-во действий, причём всякие логарифмы и функции, которые растут медленнее, чем полиномиальные, тоже входят. А неполиномиальная функция --- это экспоненциальные функция, то есть наш параметр выступает как показатель степени. Я думаю, вы сталкивались с таким алгоритмом --- любой полный перебор это всё делает. И вот иногда нам действительно нужно для поиска оптимального решения перебрать всё, абсолютно всё, а как перебрать чуть поменьше, чем всё, тема следующей лекции, а сейчас мы будем перебирать абсолютно всё. Давайте вспомним с вами основные комбинаторные объекты и формулы их количества, чтобы понимать сложность генерации всех объектов, которая может быть теоретически достижима.

\subsection*{Перестановки}
\addcontentsline{toc}{subsection}{Перестановки}
\[
    P = N!
\]

Чтобы быстро вычислять факториал, нужно раскладывать на простые множители и как-то исходя из этого ускорять вычисления. Факториал вам, скорее всего, не понадобиться, вам понадобиться остаток от деления: вам нужно будет разложить число на множители и, может быть, какое-нибудь быстрое возведение в степень вам поможет. 

\subsection*{Сочетания}
\addcontentsline{toc}{subsection}{Сочетания}
\[
    C_N^K = C_{N-1}^{K-1} + C_{N-1}^{K}
\]

Это множество без учёта порядка. То есть, мы можем говорить, что у нас есть не отсортированное $k$-элементное множество (чтобы у нас однозначно всё было), из $N$ элементов. По формуле через факториалы вычислять кол-во сочетаний мы, конечно же, не будем. Воспользуемся треугольником Паскаля, то бишь динамическим программированием. То есть сумма тех верхних стоящих элементов в треугольнике Паскаля. Ни в коем случае не через факториал!

\subsection*{Количество всех подмножеств}
\addcontentsline{toc}{subsection}{Количество всех подмножеств}
\[
    2^N
\]

Далее у нас в процессе будут возникать всякие другие штуки, например, числа Каталана.

Это всё --- неполиномиальные функции, потому что вывод будет занимать неполиномиальное время. Теперь давайте думать, как нам с этим жить. Первое, что мы должны научиться делать с комбинаторными объектами, --- это научиться их генерировать. Это можно делать разными способами. Давайте посмотрим классическую схему: она заключается в рекурсивном переборе. Начнём с простого - с подмножеств. То есть нам необходимо сгенерировать все подмножества заданного множества. Что она будет делать? Она должна будет с собой таскать уже какое-то подмножество сделанное, $N$, и позицию now. Во всех случаях, когда мы пишем рекурсию, мы проверяем позицию now: если она уже стала равна $N$, мы печатаем то самое сделанное подмножество. Во всех остальных случаях что мы делаем? Что мы делаем с очередным числом? Давайте попробуем его взять, а дальше вызовем рекурсию. Вся функция выглядит следующим образом:
\begin{verbatim}
rec(subset, N, nowpos):
    if nowpos > N:
        print(subset)
    subset.add(nowpos)
    rec(subset, N, nowpos + 1)
    subset.erase(nowpos)
    rec(subset, N, nowpos + 1)
\end{verbatim}

То есть это функция генерации всех подмножеств за $2^N$. 

Очень часто бывает так, что подмножества нужно генерировать в лексикографическом порядке. Приведенная сверху функция на такое пока не способна. Это проблему можно избежать, если вы всё это запихаете в память и отсортите, но это не круто, потому что памяти много уйдёт. Сделаем так: 
\begin{verbatim}
rec(subset, N, nownum):
    for i in nownum + 1, N
        subset.add(i)
        print(subset)
        rec(subset, N, i)
        subset.remove(i)
\end{verbatim}

Таким образом, для множества {1, 2, 3} данная функция будет генерировать следующий ответ:

1

1 2

1 2 3

1 3

2

2 3

3

Вот только пустого множества нет... Можно написать отдельную функцию, которая печатает пустую строку, или отдельно заифать, или еще что-нибудь. 

\subsection*{Генерация перестановок: по порядку и индексу}
\addcontentsline{toc}{subsection}{Генерация перестановок: по порядку и индексу}
Давайте генерить перестановки!

Как это писать? Рекурсивная функция должна постепенно накапливать перестановки, то есть в рекурсивной функции будет параметр, на какое место нам ставить элемент. И еще нам нужно каким-то образом поддерживать отсортированное множество еще неиспользованных элементов. Например, если мы уже использовали числа $3$ и $4$ при $N=5$, то множество должно быть 34125. Самое простое --- это булёвый массивчик, просто true и false. Использовали, поставили единичку и поставили false на месте неиспользованной единички. Теперь осталось только $2$ и $5$. 

То есть рекурсивная функция принимает в качестве параметров текущую позицию, некий массив, где хранится перестановка, и множество неиспользованных переменных. На каждом шаге рекурсивной функции мы проверяем, что мы не вышли за пределы $N$. После того, как вы напечатали на эту позицию, вы поочередно подставляете все числа из отсортированного множества. То есть вы идете фориком по всем числам от $1$ до $N$. Если для этого числа в массивчике стоит false, т.е. оно неиспользованное, то вы его ставите на эту позицию, помечаете как использованное, вызываете рекурсивную функцию для pos + 1, и она всеми возможными способами расставляет оставшиеся числа. Мы на каждом шаге знаем позицию и на эту позицию пытаемся расставить все числа из упорядоченного массива. Просто фориком идёт, если число не использовано еще, то ставим его, помечаем как использованное и делаем рекурсию. Сложность этого алгоритма чуть больше, чем $N!$. 

Давайте выпишем все перестановки длинной три и мы научимся всякие полезные вещи делать:

123

132

213

231

312

321

\begin{figure}[h]
\centering
\includegraphics[width=0.64\linewidth]{gust1.png}
\end{figure}

В С++ есть стандартная операция next permutation. Она получает вектор, два оператора и делает следующую перестановку. Она еще возвращает true или false: true, если всё хорошо, и false, если это была последняя перестановка и всё сломалось. То есть, вы можете написать такую штуку:

\begin{verbatim}
do {
...
...
...
} while (np(v.b()), v.e())
\end{verbatim}

И таким образом что-то там делать со всеми возможными перестановками. То есть положить числа от $1$ до $N$ и do-while сделать. Как только вы обработали последнюю перестановку, у вас функция вернёт next permutation false и ваш цикл закончится. Теперь напишем свой next permutation. Пусть у нас есть некоторая перестановка: мы смотрим на неё справа налево. У нас до какого-то момента идёт возрастание справа налево. Когда возрастание заканчивается, мы этот хвост разворачиваем полностью. Мы имеем: совпадающий у двух перестановок префикс, затем число $x$, после которого идёт возрастание справа налево, и некий суффикс, который идёт за $x$, причём утверждается, что у суффикса нет следующей перестановки. Дойдя до $x$, оно замениться на число, большее чем $x$ (вы ведь помните определение лексикографического порядка?). На какое число больше, чем $x$, оно замениться? Оно замениться \textbf{на первое число в суффиксе}, которое больше, чем $x$. Подытожим: чтобы перестановка была следующей, нужно заменить $x$ на число большее, чем $x$, а чтобы получить перестановку \textbf{ровно} следующую, нужно взять минимальное число, большее чем $x$. Со второй перестановкой понятно: первая уже больше, потому $x$ больше, а дальше суффикс сортить по возрастанию. Это сделать не трудно, потому что он у вас и так отсорчен, только в другую сторону. 

Повторим: чтобы получить из одной перестановки следующую перестановку, мы идём справо налево до тех пор, пока идёт возрастание. Находим первое число, которое выбилось из возрастания (это $x$). Далее, в этом хвосте, мы ищем число, больше чем $x$, тоже линейным поиском. Swapаем два этих числа, и теперь утверждается, что мы можем суффикс задом наперед выписать.
\begin{figure}[h]
\centering
\includegraphics[width=0.64\linewidth]{gust2.png}
\end{figure}
\newpage
Теперь мы с вами научимся генерации объекта по номеру и номера по объекту. Давайте снова рассмотрим перестановки: определим номер некой перестановки, состоящей из четырёх элементов: 4132. Если первая четвёрка, то это значит, что перед нашей перестановкой, стояли все перестановки с первой единичкой, с первой тройкой и с первой двойкой. Сколько таких было?

$4 - 1 = 3$ штуки

А сколько у нас перестановок с первой единичкой, с первой тройкой, с первой двойкой? 

$3 \times 3!$ --- для трёх штук мы переставляем три оставшиеся. Это кол-во перестановок с первой единичкой, тройкой и двойкой. Значит, номер нашей перестановки точно не меньше чем $3 \times 3!$. Окей, переходим к следующей цифре: а сколько перестановок здесь мы пропускаем? То есть сколько перестановок, где первая цифра $4$, а вторая --- меньше единицы? Да никакая. Пошла третья цифра. Сколько перестановок, где сначала 41, а третья цифра меньше трёх? Хочется сказать, что меньше трёх две цифры, это 1 и 2, но мы единичку уже брали, поэтому нам надо взять количество еще неиспользованных чисел меньше чем наше, то есть только двойку: $3 \times 3! + 1 \times 1!$.\footnote{Спасибо Полине Феофановой за найденную ошибку!} А с последней цифрой ничего не делаем, ведь она определяется однозначно.

Теперь решим обратную задачу. Допустим, нам нужно сгенерить 20-ую при нумерации с нуля перестановку из четырёх элементов. Мы по номеру будет генерить комбинаторный объект. Какая первая цифры будет в этой перестановке? Мы знаем, что для каждой первой цифры существует $3!$ перестановок, поскольку $N = 4$. Соответственно, мы делим $20$ на $3!=6$. Получаем $3$. Это означает, что у нас полностью обработано три блока с первой цифрой и сейчас идёт четвёртый. Значит, первая цифра $4$. Фиксируем, а теперь решаем задачу меньшей размерности. У нас $18$ перестановок обработано: $6$ штук с первой единичкой, $6$ штук с первой двойкой и $6$ штук с первой тройкой. Поняли, что первая цифра $4$. Теперь у нас номер перестановки равен $20-18=2$ и $N=3$. Мы делим номер перестановки на $2!$. Значит, выписываем цифру $2$, потому что первый блок уже обработан. Получили номер $0$ и $N=2$. Сколько у нас существует перестановок из 2 элементов для каждой цифры? 1! Делим $0$ на $1!$, получаем $0$ и прибавляем кол-во полных блоков (ноль), ставим единичку и выписываем это число. После этого $N=1$ и номер перестановки становится равен $0$. Итоговый ответ это 4211. Но есть проблема: это не перестановка. Что делать? Нам нужно выписывать не просто $k$-ую цифру, а $k$-ую \textbf{неиспользованную} цифру, где $k$ --- это как раз номер перестановки, который у нас есть, делённый на $(N-1)!$, то есть количество полных блоков с каждой цифры. То есть ответ на самом деле 4213. 

\subsection*{Генерация сочетаний}
\addcontentsline{toc}{subsection}{Генерация сочетаний}
Теперь научимся генерить $C_N^K$. Это как подмножество, только вы еще должны хранить кол-во элементов, которые вы должны собрать. Если вы уже всё собрали, то уже не берете тогда уже нечего. 

\subsection*{Разбиение множества на все возможные подмножества}
\addcontentsline{toc}{subsection}{Разбиение множества на все возможные подмножества}
Теперь рассмотрим такой непопулярный комбинаторный объект, как разбиение множества на все возможные подмножества. Мы будем идти рекурсивной функцией глубинной $N$ и на каждом шаге мы принимаем решение: либо кладём число уже в одно из существующих множеств, которые были созданы до этого момента, либо он создаёт новое множество и ложиться в него. Как это будет выглядеть внутри рекурсивной функции? У нас есть параметр, хранящий кол-во подмножеств, которые на данный момент у нас имеются, типа cnt subsets и очередной элемент с номером now. Массив мы назовём subset part и мы просто пройдёмся фориком от $1$ до cnt subsets и очередной элемент положим в это множество. И вызовем рекурсивную функцию. Мы дошли до элемента now и мы знаем кол-во подмножеств, которые были до этого. И мы пытаемся наш элемент положить в каждое из предыдущих множеств. Чтобы создать новое множество, вне цикла прописываем добавление в него элемента. 
\begin{verbatim}
for i in 1, cnt...
    subset part[now] = i
    rec(cnt..., now + 1)
subset part[now] = cnt + 1
rec(cnt + 1, now + 1)
\end{verbatim}

\begin{figure}[h]
\centering
\includegraphics[width=0.5\linewidth]{gust3.png}
\end{figure}

\paragraph{Задача про пиратов и ёжиков.}
Корабли противника имеют свои координаты, вектор скорости и двигаются как-то в море с постоянной скоростью. У нас есть набор атомных бомб с радиусом действия $r$. И вот нам нужно уничтожить корабли противника наименьшим кол-вом атомных бомб. Давайте пока корабли стоят на месте. Получается, нам нужно разбить корабли на все возможные подмножества и смотреть, а можем ли мы снести это подмножество одной бомбой или нет? Если нет, то отлично, просто найдём такое разбиение, где самое маленькое кол-во подмножеств. Как найти? Ну-у-у, радиус описанной окружности посчитать. Надо построить треугольник и вписать в него окружность. Теперь резко вспоминаем, что корабли противника еще и плавают. Ну или представьте, что это ёжики, и мы их должны одним ударом накормить, так сказать. Кхм. Вот смотрите, они ползают каждый под своим вектором. Как в процессе меняется радиус описанной окружности, который их всех охватывает? Такая функция будет сначала убывать, а потом возрастать. Минимум функции мы найдём бинпоиском (троичным поиском?). То есть мы строим ось $t$ и находим момент времени, когда радиус этой окружности минимален и сравниваем с радиусом нашей благодати, которую мы можем на них сослать. Если меньше или равно, то всё хорошо, если больше, то всё плохо, одной бомбы нам не хватит (в смысле крошек, мы же ёжиков кормим...). Вот такая задача, где практически используется разбиение множества на подмножества. 
\newpage
\begin{figure}[h]
\centering
\includegraphics[width=0.64\linewidth]{gust4.png}
\end{figure}

























\newpage

\section*{Лекция 14}
\section*{Перебор}
\addcontentsline{toc}{section}{Лекция 14: Перебор}
Сегодня мы изучаем метод оптимизации перебора. Будем разбирать эту тему на примере задачи о коробейнике. Коробейник --- это древний человек, который доставляет заказы, которые вы не заказывали. 
\paragraph{Задача о коробейнике.} 
У нас есть граф, для простоты пусть будет полный. И нам необходимо от начальной вершины посетить все, а дальше либо вернуться в первую, либо закончить. Вы можете это представить себе как развоз каких-то заказов. А сделать это нужно за наименьшее время. Это называется Гамильтонов цикл --- цикл, проходящий по всем вершинам. В полной графе такой точно будет, а если граф неполный, то вы можете достроить его до полного, посчитав кратчайшие пути между вершинами через промежуточные, а можете и не достраивать: это принципиально ничего не изменить. Понятное дело, в реальной жизни можно сделать задачу с полным графом. Жадный алгоритм тут работать не будет. Это задача принадлежит к классу NP. Придумаем самое простое решение: прогенерим всевозможные перестановки, в каком порядке мы эти пункты выдачи посещаем. Просто для каждой перестановки считать расстояние, которое наша газель проедет. Условно, если мы пишем рекурсивную функцию, дерево перебора будет выглядеть так: на первом шаге для $N=4$, на втором шаге на один меньше, то есть $N=3$ и т.д.

См. картинку:

\begin{figure}[h]
\centering
\includegraphics[width=0.64\linewidth]{gust5.png}
\caption{Тут изображены: полный граф, контрпример для жадного алгоритма и дерево перебора}
\end{figure}

В случае, если мы генерим все перестановки и в конце считаем маршрут, то в конце получаем, что у нас там получилось. Из всех вариантов, из всех листьев в этом дереве мы выбираем лучший. \textbf{Оценка сложности: $O(N!)$}. Это очень медленно. Если бы мы сразу в какой-то момент отсекали целое поддерево. Мы можем это сделать, когда знаем, что тут заведомо нет правильного ответа. Как это понять? Мы будем считать не в конце пройденный маршрут, а каждый раз, когда едем куда-то, по ребру конкретному. Вот мы выбрали начальную вершину, выбрали вторую и сразу прибавили к этому маршруту вес ребра между ними. Таким образом мы уже поддерживаем ту дистанцию, которая газелька проехала. И если газелька уже проехала больше, чем найденный оптимальный вариант, то ничего хорошего нас не ждет. То есть мы поддерживаем пройденный путь и лучший найденный ответ. В начале лучший найденный ответ это $\infty$. Всё понятно? Вырубаем, если нашли лучше. А как получить самое лучшее решение в начале? Запустить жадный алгоритм! Во многих случаях это может быть правильным ответом, и многие поддеревья будут сразу отрубаться, и наш перебор будет гораздо меньше. 

Есть еще один способ оптимизации. Когда мы находимся в некой вершине нашего дерева перебора, мы можем прикинуть, а сколько нам вообще ехать-то осталось. Вы находите в графе самое маленькое ребро, смотрите, сколько вообще вершин вам осталось посетить, и умножаете их друг на друга. Это оценка снизу на оставшийся путь. Вы должны суммировать то, что вы проехали, плюс оценка оставшегося пути. В нашем случае это минимальное ребро умножить на кол-во оставшихся вершин. И вот если вся эта сумма оказалась больше, чем найденный результат, то отрубаем поддерево. Нас интересует только те функции оценки, которые работают за $O(1)$. Мы уже придумали свою функцию оценки, но давайте придумаем еще лучше. Поскольку у нас Гамильтонов цикл, мы должны из каждой вершины приехать и из каждой вершины уехать. Давайте из каждой вершины уезжать по минимальному ребру, который из нее выходит. Мы вначале считаем сумму, как только посещаем какую-то вершину, вычитаем из нее вес минимального ребра и получаем сумма всех ещё не посещенных. Когда рекурсия откатывается, мы возвращаем к оценке, добавляем это число. Из оценки мы вычитаем минимальное, а к пройденному пути прибавляем то, которое мы реально прошли. И во всех графах, кроме тех, где все веса одинаковые, мы получим более хорошую оценку, чем в предыдущем нашем алгоритме.

Но этого по-прежнему мало, поэтому воспользуемся динамическим программированием. Пусть есть множество вершин, и мы посетили уже некоторое подмножество вершин subset и сейчас находимся в вершине last. Перебор будет ходить по subset`у кучей разных способов, и многие из них будут не будут отсекаться, пока subset маленький. И для каждого из них вы будете выезжать из last и дальше что-то там ездить. Понимаете, да? Если мы посетили подмножество пунктов выдачи и оказались в вершине last, нас интересует только самый короткий путь, который позволяет нам оказаться в этом состоянии. Будет решать этот вопрос двумерной динамикой: $d[subset][last]$. Чтобы попасть в last, нам нужно пройти subset/{last}.

\begin{verbatim}
d[subset][last] = min(d[subset/{last}][prev] + g[prev][last])
\end{verbatim}
где $prev \in  subset/\{last\}$

Но остались еще вопросы по поводу деталей. Поскольку у нас $N$ вершин, у нас $2^N$ subset`ов. Для каждого subset`а мы должны посчитать $N$ вариантов последней вершины. И того мы тратим $2^N \times N$ на хранение и по времени $2^N \times N:2$. Если $N$ у нас до $30$, то это очень хорошо, потому что тогда мы можем хранить subset как число. 
\paragraph{Задача о битах.} 
Необходимо установить $i$-ый бит. Если я хочу в переменную x записать $i$-ый бит, я запишу 
\\ \\
x |= 1 < < i
\\ \\
Как обнулить этот бит? Нам пригодится такая операция, когда из subset`а будет исключать last. 
\\ \\
x \&= $\sim$ (1 < < i)
\\ \\
Посмотреть на бит:
\\ \\
x \& (1 < < i)
\\ \\
Это и есть операции для множества «добавь», «удали» и «проверь». \\
\\

Теперь база динамики: мы начинаем с того, что первая вершина это $1$, а пройденное расстояние --- это $0$. Затем порядок обхода. 

\begin{figure}[h]
\centering
\includegraphics[width=1\linewidth]{gust6.png}
\end{figure}

Мы будем идти в порядке возрастания чисел
\begin{verbatim}
for subset = 1, 2^N - 1
    for last = 1, N
        if last in subset <-- здесь битовая операция
            for prev = 1, N
                if prev != last &&prev in subset
                    d[subset][last] = min(d[subset/{last}][prev] + g[prev][last], 
                                          d[subset][last}]);
\end{verbatim}

Вот код: я перебираю subset`ы в порядке возрастания, перебираю последнюю вершину, если последняя вершина лежит в подмножестве, то я перебираю предпоследнюю вершину, если предпоследняя вершина не равна последней, то применяем чудесную формулу. Насчёт инициализации: всё бесконечность, кроме первой и последней вершины --- они равны единички. Работает этот код за $2^N\times N^2$.

\paragraph{Задача о расстановке ферзей.} 
На шахматном поле размера $N\times N$ надо расставить ферзей так, чтобы они не били друг друга. Надо посчитать кол-во способов расставить их так, чтобы они друг друга не били. Мы можем решить эту задачу на $N!\times N^2$: в каждом столбце будет записан номер строки, в котором стоит ферзь, и они все должны быть различны. Решаем с помощью дерево перебора, а также делаем проверку. Как это оптимизировать? Как и в прошлой задаче: мы генерим перестановку, и каждый раз, когда поставим нового ферзя, и смотрим, нет ли на каких-то диагоналях слева от нас уже стоят ферзи. Чтобы делать это за $O(1)$, диагонали можно нумеровать, где каждая диагональ будет обладать индивидуальной числовой характеристикой. То есть делаем булёвый массивчик диагоналей, и когда у нас очередной ферзь ставится, мы ставим две пометочки на диагоналях, точка пересечения которых --- это расположение ферзя. Соответственно, нам нельзя ставить нового ферзя туда, где уже занято. Теперь мы можем это ускорить еще в два раза: на доске четного размера это наглядно видно. Я первого ферзя буду пытаться ставить только в верхнюю часть доски, и посчитать кол-во расстановок ферзей, когда первый стоит в верхней половине. И утверждается, что правильных ответов вдвое большое, потому что я могу это всё отразить во вторую половину. Для нечетного размера нам придётся считать еще и для средней линии отдельно. Ускорим этот код еще больше: в динамике у нас будут выбитые горизонтали и диагонали, то есть нам не придётся ее хранить. Это будет, по-видимому, какой-то хэш, то есть множество или словарь... Обход тут совершенно не понятен. Поэтому мы будем делать что-то типа обхода в ширину по состояниям, что мы находимся в некотором состоянии, что выбиты некоторые горизонтали и диагонали, в которые мы можем шагнуть. Таким образом, мы получим очередь новых состояний, которые добавим в очередь, и так будет работать. Здесь выгода крайне очевидна, ведь повторы будут возникать крайне редко. Но лучше пользоваться деревом, ведь оптимизация динамикой кране сложная. В домашних задачах может пригодиться такая мысль: \textbf{до какого-то момента набирать жадно, а когда осталось чуть-чуть, запустить перебор}. 





















\newpage

\section*{Лекция 15}
\section*{Корневая оптимизация}
\addcontentsline{toc}{section}{Лекция 15: Корневая оптимизация}
\subsection*{Массив префиксных сумм}
\addcontentsline{toc}{subsection}{Массив префиксных сумм}
Сегодня мы научимся делать что-то полезное с помощью выражения $\sqrt{x} = \frac{x}{\sqrt{x}}$. Например, мы забыли префиксные суммы и хотим научиться считать сумму на отрезке как-нибудь за линейное время. Как это можно сделать? Мы можем взять весь наш массив и разбить на одинаковые части, каждая из которых будет размером примерно $\sqrt{N}$. Соответственно, кол-во таких частей тоже будет $\sqrt{N}$. И мы заранее можем на каждом таком кусочке посчитать сумму. То есть мы нарезали, и сумму элементов каждого кусочка сложили в отдельный массивчик длинной $\sqrt{N}$. Теперь нам приходит запрос, мы смотрим, какие отрезки попали полностью, и возьмём эти значения из массивчика, элементы которого --- это суммы. Это займёт не более, чем $\sqrt{N}$. Но у нас могут остаться хвосты слева и справа. Эти хвостики точно меньше размера одного целого блока. То есть хвосты имеют размер $< \sqrt{N}$. То есть на запрос мы ответили за $O(3\sqrt{N})$ в худшем случае. Конечно, вы знаете способы лучше, например, за $O(1)$ или $O(log N)$ через дерево отрезков. И если у вас есть операция изменения одного элемента, то у вас на префиксных суммах этого сделать нельзя. А здесь можно! Вы можете поменять один элемент, и в нашем массивчике сумм точно так же поменять значение суммы. Если вы меняете $a$ на $b$, то в ячейку с суммой этого отрезка, где лежало $a$, добавляете разность $b-a$. 

\begin{figure}[h]
\centering
\includegraphics[width=0.5\linewidth]{gust1001.png}
\end{figure}




\newpage
Во многих случаях такой алгоритм применим, а константа у него очень маленькая. Это можно написать еще эффективнее: вместо размера блоков $\sqrt{N}$ делать их размером степени двойки. 

\subsection*{Изменения в массиве префиксных сумм}
\addcontentsline{toc}{subsection}{Изменения в массиве префиксных сумм}
Надеюсь вы помните, что префиксные суммы умеют отвечать на запрос суммы за $O(1)$. Раньше мы не могли поменять элемент, потому что иначе менялся почти весь массив префиксных сумм. Теперь мы избежим этой проблемы: мы будет запросы и изменения накапливать. Мы будет накапливать их в буфер, размер которого не будет превосходить $\sqrt{N}$ ($\leq\sqrt{N}$) Что происходит? У нас есть еще не применённые запросы (изменения). Когда буфер наполниться, мы просто возьмём и применим их \textbf{все}, а потом полностью перестраиваем массив префиксных сумм. Сложность всего этого --- это $O\left(\frac{Q}{\sqrt{N}} \times N\right) = O(Q \sqrt{N})$, где $Q$ --- кол-во запросов, $\frac{Q}{\sqrt{N}}$ --- кол-во заполнений буфера и $N$ --- перестройка. Теперь нам приходит запрос суммы на отрезке. Мы как обычно считаем сумму на одном префиксе, на другом, а далее делаем вычитание одного из другого. Но затем мы должны посмотреть в буфере запросов, какие из запросов попали в этот отрезок, и соответствующим образом изменить сумму. Такой алгоритм работает за $O(Q \sqrt{N})$.

\subsection*{Алгоритм Мо}
\addcontentsline{toc}{subsection}{Алгоритм Мо}
Данный алгоритм Мо применим только для офлайн задач! Мы опять забыли, как считается префиксная сумма... Нам даны все запросы и нам в конце нужно ответить на все запросы. Мы отсортим отрезки по их правому концу и разложим их по корзинкам, размер которых $\sqrt{N}$ каждая, в количестве $\sqrt{N}$. То есть мы можем использовать некий аналог сортировки подсчётом. Мы разложили отрезки руководствуясь нахождению их правых отрезков. Внутри последней корзинки мы отсортим отрезки по их левых концам (см. рис. на след. страничке). Потом мы смотрим на самую правую корзинку. Все правые концы отрезков лежат тут. И внутри этой корзинки мы отсортили наши отрезки по убыванию их левого конца. На всё это мы потратим $NlogN$ времени в худшем случае, если все отрезки попали в одну корзинку. Теперь поехали.

Для первого отрезка мы тупо считаем циклом for для суммы за $N$. 

\begin{figure}[h]
\centering
\includegraphics[width=0.5\linewidth]{gust1002.png}
\caption{Отрезки и корзиночки}
\end{figure}

Переходим ко второму отрезку: утверждается, что между концами второго и первого отрезка разность не более $\sqrt{N}$. А дальше мы просто будем считать сумму только того торчащего кусочка, край которого находится левее левого края первого отрезка. Затем третий отрезок. Убрали из суммы второго отрезка те элементы, которые \textbf{не} входят в третий. И только потом мы прибавляем к этой сумме те элементы, которые лежат в третьем отрезке и не лежат ни в каком другом, ведь левый край третьего отрезка самый левый из всех краёв (ведь они посорчены).

\begin{figure}[h]
\centering
\includegraphics[width=0.5\linewidth]{gust1003.png}
\caption{Отрезки с учетом всех действий}
\end{figure}

Сложность всего этого $O(Q\sqrt{N} + N\sqrt{N}) = O(N \sqrt{N})$, где $Q=N$, но не всегда. 

\paragraph{Задача о поиске медианы.} 
Вы делаете два set`а (в одном маленькие элементы, а в другом большие) и, когда ваш отрезок сдвигается, один элемент, который выпал, убираете из set`ов, а тот элемент, который попал, наоборот добавляете, затем выравниваете их размеры, и самый маленький элемент второго set`a будет медианой. 

\paragraph{Задача о множестве строк.} 
Пусть у нас есть множество строк и операции <<добавь>> и <<удалить>>. И еще у нас есть операция <<определи, сколько раз подстрока входит во все строки текущего множества>>. В таких задачах часто бывает ограничение на суммарную длину строк в запросах. То есть длина входных данных ограничена $l$. Назовём строки \textit{лёгкими}, если их длина не превосходит $\sqrt{l}$, и \textit{тяжёлыми} иначе. Но кол-во тяжёлых $\leq \sqrt{l}$. Допустим, к нам пришла строка, размер которой меньше $l$. Мы переберём вообще все ее подстроки, какие есть (с помощью двух циклов), чтобы посчитать их хэши. А потом создаём словарь, где ключём является хэш, а значением сколько раз эта подстрока встречалась. А когда придёт запрос на одну из этих подстрок, мы посчитаем его за длину запроса. Это для лёгких строк, а что делать с тяжёлыми? Пусть строки тяжёлые, а запросы лёгкие. Мы переберём в них все подстроки, длина которых не превосходит $\sqrt{l}$. А дальше точно так же: ищем хэш, кладём в словарь.

\begin{figure}[h]
\centering
\includegraphics[width=0.5\linewidth]{gust1004.png}
\caption{Кол-во подстрок и их название}
\end{figure}

Осталось самое последнее: решить что делать с тяжёлыми запросами. Кол-во таких запросов $\leq \sqrt{l}$. Тяжёлые запросы мы будем искать только в тяжёлых строках. Как? Хэшами. Вот это окошко двигать и всё. Общее время ответа на все тяжёлые запросы $l \sqrt{l}$. Таким образом, вся задачка решилась за $l \sqrt{l}$


\paragraph{Задача о треугольниках графа.} 
Дан неориентированный граф, где $n = m$ (плотных графов тут, скорее всего, не будет). Необходимо найти кол-во треугольников в нём. Ограничения: $n$ и $m$. Давайте разделим наши вершины на лёгкие и тяжёлые. Лёгкие --- это те, у которых $\leq \sqrt{N}$ соседей, а тяжёлые иначе. 

\renewcommand{\arraystretch}{1.7} %% increase table row spacing
\renewcommand{\tabcolsep}{0.25cm}   %% increase table column spacing
\begin{center}
\centering
\begin{tabular}{|c|c|c|}
\hline
--- & Соседей & Количество \\
\hline
Лёгкие & $\leq \sqrt{N}$ & $\leq N$ \\ 
\hline
Тяжёлые & $> \sqrt{N}$ & $\leq \frac{M}{\sqrt{N}}$\\
\hline
\end{tabular}
\end{center}

\textbf{Случай 1}: все вершины лёгкие. Точнее зафиксируем ребро $ab$, и оба его конца оказались лёгкими. Если третья вершина лёгкая, как нам посчитать кол-во таких треугольников? Посмотрели всех соседей и пересекли это множество. Сложность этого $O(M \sqrt{N}$. 

\textbf{Случай 2}: одна вершина лёгкая, а другая тяжёлая. Нам нужно посчитать кол-во её соседней. Мы можем взять всех соседей лёгкой вершины и в множестве за $O(1)$ проверить, есть ли они в списке соседей тяжёлой вершины. Сложность этого $O(M \sqrt{N})$.

\textbf{Случай 3}: обе вершины тяжёлые.

\noindent \hspace{4em} \textbf{Подслучай 3.1}: обе тяжёлые и одна лёгкая. Если считать аккуратно, то они у нас зачтутся в первых двух случаях. 

\noindent \hspace{4em} \textbf{Подслучай 3.2}: обе тяжёлые и третья тяжёлая. У нас есть $N \sqrt{N}$ способов выбрать 3 тяжёлых вершины. Мы храним множество тяжёлых соседей дополнительно у каждой тяжёлой вершины. 
\begin{figure}[h]
\centering
\includegraphics[width=0.52\linewidth]{gust1005.png}
\caption{Тем временем на доске...}
\end{figure}

Пишите аккуратно, чтобы не зачесть один и тот же треугольник по несколько раз!

\subsection*{Суммы и сдвиги}
\addcontentsline{toc}{subsection}{Суммы и сдвиги}
Нужно найти сумму на отрезке от $L$ до $R$ и сделать циклический сдвиг этого кусочка. Смотрите, если мы нарезали всё на блоки длины $\sqrt{N}$, то когда у нас происходит какое-то изменение на отрезке, то как пересчитываются суммы? Как раньше делать не получится: выходит какая-то жесть. Давайте посмотрим, что происходит при циклическом сдвиге на единичку влево. Храниться они будут, например, в виде дека. Мы элемент из четвертого дека перекладываем в третий в конец, из третьего во второй в конец и так далее (см. картинку). А в конце берем элемент из серединки первого и кладём в серединку четвертого. $\sqrt{N}$ операций. 

\begin{figure}[h]
\centering
\includegraphics[width=0.52\linewidth]{gust1006.png}
\end{figure}



\paragraph{Задача о смене местожительства.}
Теперь у нас блоки переменного размера! Нам нужно i-ый элемент переставить на другое место. Сначала всё режем на $\sqrt{N}$. Когда нам приходит запрос, то мы будем перестраивать каждый корень из $N$ запросов, полностью заново обновлять (идея с буферами). Но наша логика может порушиться, если бездумно так поступать. Поэтому как только блок превысит размер $2 \sqrt{N}$ (я это может случиться не раньше чем раз в $\sqrt{N}$ запросов) мы просто возьмём и всё перестроим заново за линейное время $O(N)$, то есть нарежем аккуратненько на блоки одинакового размера. Это может случиться $\sqrt{Q}$ раз и занимает время $N$. Так же, нам понадобиться одномерный массив, в котором будут храниться размеры блоков. По мере возникновения операций его нужно обновлять. 

\newpage
\begin{figure}[h]
\centering
\includegraphics[width=1\linewidth]{gust1007.png}
\caption{Тем временем на доске...}
\end{figure}



\newpage 
Выражаю благодарность своим самым преданным читателям!
\begin{figure}[h]
\centering

\includegraphics[width=0.2\linewidth]{gust99999.jpg}
\end{figure}





\begin{center}
\fontsize{70}{80}\selectfont\textbf{КОНЕЦ}
\end{center}


\end{document}